% 图 2-Z : 窗口内容震荡与节点指纹匹配原理图
\begin{figure}[htbp]
\centering
\begin{tikzpicture}[
    node distance=2cm,
    block/.style={draw, rectangle, rounded corners, minimum width=2.5cm, minimum height=1cm, text centered, font=\sffamily\small, fill=blue!5},
    ghost/.style={draw, rectangle, dashed, minimum width=2.5cm, minimum height=1cm, text centered, font=\sffamily\small, fill=gray!20, text=gray},
    arrow/.style={thick,->,>=stealth},
    label/.style={font=\sffamily\scriptsize, text=red!70!black}
]

% 左侧：黑盒化混淆状态
\node[font=\sffamily\bfseries, text=black] (title1) {步骤一：混淆隔离态 (Ghost Nodes)};
\node[ghost, below of=title1, yshift=1cm] (g1) {透明占位节点};
\node[ghost, below of=g1, yshift=0.5cm] (g2) {混淆资源 ID};
\node[ghost, below of=g2, yshift=0.5cm] (g3) {隐藏真实控件};

% 中间：主动震荡探测
\node[block, fill=orange!20, right of=g2, xshift=2.5cm] (shake) {窗口内容震荡\\(微位移滑动 1-2px)};
\draw[arrow, red, ultra thick, dashed] (g2.east) -- node[above, font=\sffamily\scriptsize, text=red] {触发更新} (shake.west);

% 右侧：真实树挂载与匹配
\node[font=\sffamily\bfseries, text=black, right of=title1, xshift=7cm] (title2) {步骤二：真实树挂载与特征定位};
\node[block, fill=green!10, below of=title2, yshift=1cm] (r1) {真实控件子树暴露};
\node[block, fill=green!10, below of=r1, yshift=0.5cm] (r2) {拓扑深度提取};
\node[block, fill=green!10, below of=r2, yshift=0.5cm] (r3) {计算视觉包围盒重心};

\draw[arrow, blue, thick] (shake.east) -- node[above, font=\sffamily\scriptsize, text=blue] {强制重绘 (Remount)} (r2.west);

% 下方：特征指纹匹配模型
\node[draw, rectangle, minimum width=10cm, minimum height=1.5cm, fill=yellow!10, below of=r3, yshift=-0.5cm, xshift=-2.5cm] (fingerprint) {};
\node[font=\sffamily\bfseries, text=black, xshift=-2cm, yshift=0.4cm] at (fingerprint.center) {UI 特征归一化模型 (Feature Normalization)};
\node[font=\sffamily\scriptsize, text=black, yshift=-0.2cm] at (fingerprint.center) {匹配度 $P = w_1 \times (相对拓扑关系) + w_2 \times (视觉重心偏移量) > 98\% \Rightarrow$ \textbf{锁定目标}};

\draw[arrow, thick] (r3.south) -- (fingerprint.north -| r3.south);

\end{tikzpicture}
\caption{基于窗口内容震荡与层级拓扑指纹的主动探测匹配机制}
\label{fig:wechat_shake_match}
\end{figure}

% -------------------------------------------------------------

% 图 2-W : 焦点退避与心跳防死锁机制时序图
\begin{figure}[htbp]
\centering
\begin{tikzpicture}[
    node distance=2cm,
    actor/.style={draw, rectangle, minimum width=2.5cm, minimum height=0.8cm, text centered, font=\sffamily\bfseries, fill=blue!10},
    timeline/.style={thick, dashed},
    arrow/.style={thick,->,>=stealth},
    message/.style={font=\sffamily\scriptsize, text=black}
]

% 角色定义
\node[actor] (user) {老年用户 (干扰源)};
\node[actor, right of=user, xshift=2cm] (desktop) {桌面主程序};
\node[actor, right of=desktop, xshift=3cm] (accessibility) {无障碍服务引擎};

% 时间线
\draw[timeline] (user.south) -- +(0,-5.5);
\draw[timeline] (desktop.south) -- +(0,-5.5);
\draw[timeline] (accessibility.south) -- +(0,-5.5);

% 步骤 1: 触发与焦点退避
\draw[arrow] ([yshift=-1cm]desktop.south) -- node[above, message] {1. 自动化任务发起} ([yshift=-1cm]accessibility.south);
\node[draw, fill=orange!20, font=\sffamily\tiny, right of=desktop, xshift=1cm, yshift=-1.5cm] {2. 释放焦点，隐藏悬浮窗};
\draw[arrow, ->, bend right] ([yshift=-1cm]desktop.south) to ([yshift=-2cm]desktop.south);

% 步骤 2: 用户误触干扰
\draw[arrow, red, thick] ([yshift=-2.5cm]user.south) -- node[above, message, text=red] {3. 误触屏幕 / 权限弹窗} ([yshift=-2.5cm]accessibility.south);

% 步骤 3: 防死锁心跳检测
\node[draw, fill=red!10, font=\sffamily\tiny, right of=accessibility, xshift=1.5cm, yshift=-3.5cm] {4. 状态机停滞 > 5秒};
\draw[arrow, ->, bend left, thick, dashed] ([yshift=-3cm]accessibility.south) to ([yshift=-4cm]accessibility.south);
\node[message, right of=accessibility, xshift=1.5cm, yshift=-4.2cm] {触发死锁判定};

% 步骤 4: 自回归与恢复
\draw[arrow, blue, thick] ([yshift=-5cm]accessibility.south) -- node[above, message, text=blue] {5. 强制模拟返回 (Self-regression)} ([yshift=-5cm]desktop.south);

\end{tikzpicture}
\caption{面临用户误触与弹窗干扰时的焦点退避与心跳防死锁时序图}
\label{fig:focus_retreat}
\end{figure}
